\section{Results}

\subsection{Cross-dataset generalization}
Table~\ref{tab:main} \todo{generate with scripts/paper_tables.py} reports every test set in one
table: the same-dataset result sits next to the cross-dataset results, not above them.

\begin{table}[t]
\centering
\small
\caption{Video-level AUC, mean $\pm$ std over three seeds, all models trained on FF++ c23 only.}
\label{tab:main}
\begin{tabular}{lcccc}
\toprule
Test set & Ours & Plain ViT & Xception & EffNet-B4 \\
\midrule
FF++ c23 (same dataset) & n/a & n/a & n/a & n/a \\
FF++ c40 & n/a & n/a & n/a & n/a \\
Celeb-DF v2 & n/a & n/a & n/a & n/a \\
DFDC subset & n/a & n/a & n/a & n/a \\
DeeperForensics-1.0 & n/a & n/a & n/a & n/a \\
WildDeepfake & n/a & n/a & n/a & n/a \\
\midrule
Unseen FF++ method (mean) & n/a & n/a & n/a & n/a \\
ECE on Celeb-DF $\downarrow$ & n/a & n/a & n/a & n/a \\
\bottomrule
\end{tabular}

\end{table}

\todo{Two or three sentences: headline number, comparison with the strongest baseline, where
the gap closes and where it does not.}

\subsection{Unseen manipulation methods}
\todo{Table of held-out-method AUC per method and the mean (M5), and what the hardest held-out
method tells us.}

\subsection{Robustness}
\todo{Figure: AUC against JPEG quality, blur, downscaling, noise and contrast. Which degradations
hurt most, and how SBI and augmentation change the curves.}

\subsection{Calibration}
\todo{ECE before and after temperature scaling on every test set; how well the FF++ fitted bands
transfer to unseen data.}

\subsection{Fairness}
\todo{AUC and false alarm rate per demographic group where published labels exist; every gap
above 5 AUC points discussed.}

\subsection{Ablations}
\begin{table}[t]
\centering
\small
\caption{Ablations (1 seed, 8 epochs). Each row changes one setting of the final model.}
\label{tab:ablations}
\begin{tabular}{llccc}
\toprule
& Option & FF++ & Celeb-DF & Unseen \\
\midrule
\multicolumn{5}{c}{n/a: ablations not run yet} \\
\bottomrule
\end{tabular}

\end{table}

\todo{Explain why the winning settings help, with evidence, not just which ones won.}

\subsection{What the model looks at, and where it fails}
\todo{Heatmap figure on correct and wrong predictions; patterns among false alarms and misses.}
